
%!TEX TS-program = xelatex
%!TEX encoding = UTF-8 Unicode

\documentclass[english,a4paper,11pt]{article}
%\usepackage[T1]{fontenc}
%\usepackage[utf8]{inputenc}
\usepackage{}
\usepackage[margin=3cm,a4paper]{geometry}

\usepackage{graphicx}
\usepackage{amssymb}
\usepackage[norsk]{babel}
\usepackage{color}
\usepackage{xcolor}
\usepackage{booktabs}
\usepackage{array}
\usepackage{wrapfig}
\usepackage{calc}
\usepackage{hyperref}
\usepackage{xspace}
\usepackage{parskip}


\definecolor{xalblue}{rgb}{0.1992,0.1992,0.5977}
\definecolor{xaldark}{rgb}{0.1992,0.1992,0}
\definecolor{lightgray}{gray}{0.8}


\usepackage{fontspec,xltxtra,xunicode,relsize}
\defaultfontfeatures{Mapping=tex-text}
%\setmainfont[Mapping=tex-text]{Verlag Light}
%\setmainfont[Mapping=tex-text]{Sina-Regular}
%\setsansfont[Scale=MatchLowercase,Mapping=tex-text]{Helvetica Neue}
%\setmonofont[Scale=MatchLowercase]{Monaco}
\renewcommand*{\familydefault}{\sfdefault}

\newcolumntype{L}{>{\raggedright}p{0.14\textwidth}}
\newcolumntype{R}{p{0.76\textwidth}}
\newcommand\VRule{\color{lightgray}\vrule width 0.5pt}


\usepackage{titlesec}
\newcommand\secformat[1]{\normalfont\Large\bfseries\color{xalblue}#1}
\titleformat{name=\section}[block]{}{}{}{\secformat}

\begin{document}

\vspace*{-3cm}
\hspace*{12.5cm}\includegraphics[width=4cm]{logo.pdf}

\vspace*{-1.65cm}

\noindent
{\LARGE \bfseries \color{xaldark}Curriculum Vitae for Ola Skavhaug}

\section*{Personal information}

\begin{minipage}[t]{0.4\textwidth}
\begin{tabbing}Nationality---\=olaf\kill 
Address:\>Aslakveien 31, \\
\>0753 Oslo \\
Born:\>10.04.1974 \\
\end{tabbing}
\end{minipage}
$\quad\quad$
\begin{minipage}[t]{0.3\textwidth}
\begin{tabbing}Nationality---\=olaf\kill 
E-mail: \>ola@xal.no\\
Phone: \>+47 926 12 490\\
Nationality:\>Norwegian\\
\end{tabbing}
\end{minipage}

\section*{Summary}
I am a software developer, researcher, and project leader with 16 years professional experience. My main areas of technical expertise are mathematical and numerical software development, algorithm development, advanced scripting with modern scripting languages, parallel programming, software testing and deployment, library design, and scientific visualization.\\
\\
Through my work, I have obtained an extensive set of skills that allows me to understand and solve challenges in collaboration with other experts. Today's industrial challenges are often multidisciplinary, and involve competences from several fields at once. A key challenge in such projects is to communicate across disiplines.


\section*{Technical Skills}
\begin{tabular}{LR}
Languages& C, C++, C++11, Java, Fortran, Python, Javascript, Perl, PHP, Bash, Tcl/tk, Matlab, Sql, LaTeX, HTML, XML\\
Frameworks& Numpy, SciPy, Matplotlib, MPI, BSP, Swig, Boost, Stl, Armadillo, dlib, VTK, FEniCS, PETSc, SLEPc, Diffpack\\
Tools& Subversion, Mercurial, Git, cvs, Make, CMake, Scons, GCC, Autoconf, Linux, css, MySQL
\end{tabular}

\section*{Education}
\begin{tabular}{LR}
2004 & Dr. Scient in Computer Science, The Faculty of Mathematics and Natural Sciences, University of Oslo. Thesis’ title: ``Numerical Methods and Software with Applications in Computational Finance''.\\
1998 & Cand. Scient in Computer Science, Department of Informatics, University of Oslo
\end{tabular}

\section*{Professional Experience}
\begin{tabular}{LR}
2013 -- & Consultant, Expert Analytics\\
2011 -- 2013& Innovation manager at Simula Innovation\\
2010 -- 2011& Senior Scientific Programmer at Kalkulo AS\\
2007 -- 2010& Research Scientist and head of the computational middleware software activity at the Centre of Biomedical Computing (CBC) at Simula Research Laboratory\\
2005 -- 2007& Research Scientist and head of the project "Software for PDEs" at Simula Research Laboratory\\
2004 -- 2005& IT-manager, Simula Research Laboratory\\
2004& System Administrator, Simula Research Laboratory\\
2004 -- 2010& 20\% Associate Professor, Department of Informatics, University of Oslo\\
2001 -- 2004& Ph.D. student at the Simula Research Laboratory\\
2000 -- 2004& 20\% Teaching Position at the Department of Informatics, University of Oslo\\
2000 -- 2001& Ph.D. student at the Department of Informatics, University of Oslo\\
\end{tabular}

\section*{Other Experience}
\begin{tabular}{LR}
2009 -- 2013& Employee representative in the board of directors, Simula Research Laboratory\\
2005 -- 2006& Board member, Øraker Barnehage AS\\
\end{tabular}

\section*{Languages}
\begin{tabular}{LR}
Norwegian & Mother tongue \\
English & Fluent \\
German & Basic \\
\end{tabular}

\section*{Personal Skills}
\begin{tabular}{LR}
Management & Motivate and lead experts and PhD students, define and implement new projects, facilitate communication in informal surroundings to break up the work day.\\
Applied mathematics & Analyze, develop and implement complex algorithms in applied sciences, while balancing constraints like flexibility and efficiency. Short, agile development cycles with discussions and feedback from problem owners.\\
\end{tabular}

\section*{Some interests and hobbies}
\begin{tabular}{LR}
Physical & Telemark skiing, running, biking, climbing.\\
Gastronomical & Beer brewing, sausage making.\\
Other & Reading, traveling, trekking, expeditions.\\
\end{tabular}

\section*{Extended descriptions of select projects}

\begin{tabular}{@{}LR}
Activity & Operationalization of Enki Open Source for Statkraft \\
Role & Solution architect and software developer \\
Staffing & Two main developers and several researchers and students \\
Description & Enki is an open source hydrologcal simulation platform initially developed by SINTEF for Statkraft over a time period of several years. The original implementation, being a proof of concept type research  code, was not suited for daily use in an operational setting due to limitations on scalability, design and performance. I was hired in a project to remedy this, and we realized that we needed to reimplement everything from scratch and use the existing code base as a starting point for the algorithmic aspects. The resulting hydrological forcasting toolbox, named SHyFT, is now run daily at Statkraft. \\
Tools & C++11, Stl, Swig, Python, Armadillo, dlib, Boost, yaml, Codeblocks \\ 
\addlinespace \bottomrule[.1pt] \addlinespace
   \end{tabular}

\begin{tabular}{@{}LR}
Activity & POC data analysis scripting platform for Statkraft \\
Role & Solution architect and software developer \\
Staffing & Project leader, data analysts, and two main developers \\
Description & Statkraft's current analysis platform for midterm forcasting of power markets is getting old. New technical solutions that can handle both the complexity of their models in a natural way, and also changing requirements caused by massive data growth, are therefore being considered. In the project, I contributed to the implementation of a new domain model, desiphered the binary result files from the core market simulator to populate this model with data, and then made a proof of concept analysis platform in Python that showcased some of the possibilities such tool can offer. The project was considered a success, and released funding for planning a complete new analysis platform. \\
\addlinespace \bottomrule[.1pt] \addlinespace
   \end{tabular}

\begin{tabular}{@{}LR}
Activity & mCASH backend development\\
Role & Senior Python Developer\\
Staffing & 12-15 Python developers\\
Description & In this project, I have been working on most parts of the backend of a new mobile payment system. This includes financial transaction handling, the internal bank implementation, messages emitted through various protocols based on recipients, OpenID Connect scopes implementation and payment for these, web handlers for endpoints, and Datastore transaction in the Google app engine, all in Python. I have also rewritten the instrumentation test framework, that does real life scenario, black box testing of the core system. The development is test driven, with tests covering close to 100 percent of the code base, follows the Scrum agile method, and utilizes state-of-the-art technologies like Git for version control, Nose for testing, and Sphinx for documentation. \\
Tools & Python, Google app engine, Git, buildout, nose tests, Sphinx, webapp2, Jinja2, OAuthLib, JSON, html, javascript, jQuery, Pusher\\ 
\addlinespace \bottomrule[.1pt] \addlinespace
   \end{tabular}

\begin{tabular}{@{}LR}
Activity & Computational Middleware, Center for Biomedical Computing\\
Role & Leader, scientist and software developer\\
Staffing & 3-4 scientists and developers\\
Description & As a project leader, I had the responsibility of the technical development, scientific achievements, the economy, and the personnel in the project. The project strongly emphasized the development of reliable, extensible and numerical efficient software components for solving scientific problems through simulations.\\ 
\addlinespace \bottomrule[.1pt] \addlinespace
   \end{tabular}

\begin{tabular}{@{}LR}
Activity & Software for PDEs, Simula Research Laboratory\\
Role & Leader, scientist and software developer\\
Staffing & 6-8 scientists, developers and PhD students for two years\\
Description & Under my responsibility the project defined and developed novel software frameworks for advanced computer simulation and visualization based on a hybrid technology of C++ and Python. Simula Research Laboratory applied and was awarded a Centre of Excellence by the Research Council of Norway in 2007, where the activity of this project became a central component.\\ 
\addlinespace \bottomrule[.1pt] \addlinespace
   \end{tabular}

\begin{tabular}{@{}LR}
Activity & Python Computing Components\\
Role & Main developer\\
Description & PyCC is a modern and efficient scripting framework that is used to solve differential equations modelling the electrical activity in the human heart – the so called bidomain equations. The complexity of the problem, and the use of the tool to conduct research, requires both flexibility and efficiency. To meet these needs we implemented a high level scripting interface in Python for flexibility, and migrated bottlenecks to low-level extension modules implemented in C/C++ and Fortran. Central activities were library design, interface building strategies, cross language techniques, code generation, and third party software integration.\\
Tools & C/C++, Fortran, Python, Swig, MPI, Subversion, Scons, PETSc, FEniCS, Hypre, BoomerAMG, Diffpack, GNU Compiler Collection and Debugger, Valgrind\\ 
\addlinespace \bottomrule[.1pt] \addlinespace
   \end{tabular}

\begin{tabular}{@{}LR}
Activity & Viper\\
Role & Main developer\\
Description & Viper is a lightweight runtime visualization framework for scientific data and results. It grants the underlying visualization library, VTK, direct access to the simulation result, thereby minimizing memory copies for efficiency. Viper can visualize both scalar and vector data, as well as wireframe geometries (meshes).\\
Tools & C/C++, VTK, Python\\ 
\addlinespace \bottomrule[.1pt] \addlinespace
   \end{tabular}

\begin{tabular}{@{}LR}
Activity & Gotran\\
Role & Main developer\\
Description & Systems of ordinary differential equations are often complex, and implementing these in a numerically efficient way is both time consuming and error prone. To remedy this, I have implemented Gotran – a compiler that takes ODEs described in a high level DSL (domain specific language) and generates highly specific and numerically efficient C/C++ code. By building on top of another software project I have implemented, Swiginac, Gotran utilizes symbolic manipulation during several of the code transformations to reduce the number of floating point operations needed to evaluate the ODE systems during simulation.\\
Tools & C/C++, Python, Swiginac, Swig\\ 
\addlinespace \bottomrule[.1pt] \addlinespace
   \end{tabular}

\begin{tabular}{@{}LR}
Activity & Swiginac - extending Python with symbolic mathematics\\
Role & Main developer\\
Staffing & Open source project with several contributors\\
Description & Swiginac is a symbolic mathematics module for Python. It is built by exposing GiNaC, a symbolic manipulation library written in C++ to Python with Swig. The efficiency of the underlying C++ library makes Swiginac one of the fastest technologies in its class in Python, and the possibilities of writing expressions in various ways makes Swiginac well suited for code generation purposes. Swiginac was developed as a side project during my PhD, in order to make a system for automatic code verification of numerical simulators.\\
Tools & C++, Stl, Python, Swig, Distutils, Subversion, Make\\ 
\addlinespace \bottomrule[.1pt] \addlinespace
   \end{tabular}

\begin{tabular}{@{}LR}
Activity & Famms - automatic code verification for PDEs\\
Role & Main developer\\
Description & Standard PDE problems can be formulated as $F(u)=0$, where $F$ is a possibly non--linear system of differential equations. The task is then to find the unknown $u$. By selecting a manufactured, analytical solution instead, called $v$, we can compute $b=F(v)$. Then by defining $G = F(v)-b$, we again obtain a standard problem on the form $G(v)=0$. Forgetting that we know $v$, we can try to solve the last equation to see if the numerical simulator is working as expected. The method above is commonly referred to as the method of manufactures solutions, and Famms, the software system I implemented, automates this process by calculating both $b$ and the perturbed problem $G$ with minimal effort.\\
Tools & C/C++, Python, Swig, Swiginac, Diffpack, FEniCS, PyCC.\\ 
\addlinespace \bottomrule[.1pt] \addlinespace
   \end{tabular}

\begin{tabular}{@{}LR}
Activity & Biomedical computing\\
Role & Developer and project leader\\
Staffing & Two developers\\
Description & In this project, a California based software company in biomedicine wanted to incorporate some of the technology I had developed into their commercial code to strengthen the finite element analysis and visualization capabilities of their software. Over a period of six months, we successfully integrated the components into their code, such that they could use PyCC and Viper.\\ 
\addlinespace \bottomrule[.1pt] \addlinespace
   \end{tabular}

\begin{tabular}{@{}LR}
Activity & Symphonical \\
Role & Main developer \\
Staffing & Project leader and two developers \\
Description & Symphonical is a web-based collaboration tool. Initially it was conceived as a tool for running agile software development projects based on the metaphor of post-it notes on a virtual wall, a scope that since has been widened. We created the first prototype of the system, in 2005, before it was spun out as a company and developed further by others. \\
Tools & PHP, Mysql \\ 
\addlinespace \bottomrule[.1pt] \addlinespace
   \end{tabular}

\begin{tabular}{@{}LR}
Activity & Gründergarasjen, Simula Innovation\\
Role & Manager\\
Staffing & 35 entrepreneurs working on their startups\\
Description & Responsible for planning and implementing Gründergarasjen, a free startup incubator hosted by Simula. By opening up our research laboratory for entrepreneurs, we hope to provide a creative environment in which new ideas are born and realized. The garage is loosely based on the 12 rules of the HP garage; commonly recognized to be the birth place of Silicon Valley. Entrepreneurs can stay with us up to one year, provided they get positive evaluations.\\ 
\addlinespace \bottomrule[.1pt] \addlinespace
   \end{tabular}

\begin{tabular}{@{}LR}
Activity & Department of Informatics, University of Oslo\\
Role & Associate Professor\\
Description & Over a period of ten years, I have given lectures in two popular courses at the university, teaching students how to apply high-level computer languages for advanced problem solving, and how mix languages for both convenience and numerical efficiency.\\
Tools & Python, C/C++, Swig, Perl, Bash/Sh, Tcl/tk, CGI\\ 
\end{tabular}

\end{document}

